\documentclass[11pt]{article}
\usepackage[margin=1in]{geometry}
\usepackage{amsmath,amssymb,amsthm,mathtools}
\usepackage{graphicx}
\usepackage{microtype}
\usepackage[hidelinks]{hyperref}

\newtheorem{theorem}{Theorem}
\newcommand{\D}{\mathbb D}
\newcommand{\Hk}{\mathcal H(K)}
\newcommand{\Hl}{\mathcal H(L)}

\title{Vertical-Slice Reconstruction in the Canonical Agler Space\\
\large A complete RKHS answer to Question 5.6 of arXiv:1312.0177}
\author{fa\_banach\_001, lane 4}
\date{August 17, 2026}

\begin{document}
\maketitle

\begin{center}
\fbox{\parbox{0.91\textwidth}{
\textbf{Status: literature-implied answer (complete answer to Question 5.6).}
The needed mechanism is the classical restriction theorem for reproducing
kernel Hilbert spaces.  The new step here is its direct identification with
the vertical-slice uniqueness argument already proved in the source paper.
}}
\end{center}

\section{Source question}

On page 36 of \cite{BickelKnese}, Bickel and Knese ask:
\begin{quote}
\textbf{Question 5.6.} Assume \(g\in\mathcal H(K^{\max}_1)\).
Is there an explicit way to construct \(g\) using only the function
\(g(0,w_2)\)?
\end{quote}
They add that a clean answer would give formulas for two canonical
colligation blocks \(A\) and \(B\).

\begin{figure}[h]
\centering
\includegraphics[width=\textwidth]{figures/open_problem_crop.png}
\caption{Question 5.6 and its stated motivation in the source PDF, page 36.}
\end{figure}

\section{Identification and proof}

The restriction principle used below is classical; see Aronszajn's
Section 5, \emph{The restriction of a reproducing kernel} \cite{Aronszajn}.
We include the short operator-valued proof, so the conclusion does not
depend on any unstated scalar-to-vector extension.

Let \(E_*\) be the coefficient Hilbert space, let
\[
K=K^{\max}_1:\D^2\times\D^2\longrightarrow\mathcal L(E_*),
\qquad
L(\zeta,\eta):=K((0,\zeta),(0,\eta)),
\]
and let \(R:\Hk\to\Hl\) be vertical restriction,
\[
(Rg)(\zeta)=g(0,\zeta).
\]
Here \(\Hl\) means the RKHS associated with the restricted kernel \(L\),
not an ambient Hardy norm.

\begin{theorem}[Exact reconstruction from the vertical trace]
\label{thm:reconstruct}
For the canonical kernel \(K^{\max}_1\) of the source paper, \(R\) is a
unitary operator.  Hence, if \(h(\zeta)=g(0,\zeta)\), then
\[
g=R^*h.
\]
Equivalently, for every \(z\in\D^2\) and \(\nu\in E_*\),
\begin{equation}
\label{eq:cross}
\langle g(z),\nu\rangle_{E_*}
=
\left\langle
h,\ \zeta\longmapsto K((0,\zeta),z)\nu
\right\rangle_{\Hl}.
\end{equation}
Thus \(g\) is constructed using only its one-variable trace and the
known canonical kernel.
\end{theorem}

\begin{proof}
For \(\eta\in\D\) and \(u\in E_*\), write
\[
K_{(0,\eta)}u=K(\,\cdot\,,(0,\eta))u,
\qquad
L_\eta u=L(\,\cdot\,,\eta)u.
\]
On the dense span of the \(L_\eta u\), define
\[
J(L_\eta u)=K_{(0,\eta)}u.
\]
The two Gram matrices agree:
\[
\langle L_\eta u,L_\zeta v\rangle_{\Hl}
=
\langle K_{(0,\eta)}u,K_{(0,\zeta)}v\rangle_{\Hk}.
\]
Therefore \(J\) extends isometrically from \(\Hl\) into \(\Hk\).
Its range is
\[
\overline{\operatorname{span}}\{K_{(0,\eta)}u:
\eta\in\D,\ u\in E_*\}
=
\{f\in\Hk:f(0,\eta)=0\ \forall\eta\in\D\}^{\perp}.
\]
The source paper proves immediately before Question 5.6 that the vertical
slice is a uniqueness set for \(\mathcal H(K^{\max}_1)\).  The vanishing
space in the last display is consequently zero, so \(J\) is onto.

The reproducing property gives \(R=J^*\), initially on kernel functions
and hence everywhere.  Since \(J\) is unitary, so are \(R\) and \(R^*=J\),
and \(g=R^*(Rg)\).  Finally,
\[
\langle g(z),\nu\rangle
=\langle g,K_z\nu\rangle_{\Hk}
=\langle h,RK_z\nu\rangle_{\Hl},
\]
while \((RK_z\nu)(\zeta)=K((0,\zeta),z)\nu\).  This is
\eqref{eq:cross}.
\end{proof}

\section{Finite Gram reconstruction}

Formula \eqref{eq:cross} is exact, but it also gives a concrete convergent
algorithm.  Choose points \(\lambda_n\in\D\) with an accumulation point in
\(\D\), and choose a countable total set of vectors in \(E_*\).  Enumerate
the resulting kernel sections as
\[
\ell_j(\zeta)=L(\zeta,\lambda_j)u_j,\qquad j\geq1,
\]
repeating points as needed to pair them with coefficient vectors.  Let
\(P_N\) denote the orthogonal projection in \(\Hl\) onto
\(\operatorname{span}\{\ell_1,\ldots,\ell_N\}\).

The coefficients \(c^{(N)}_j\) of \(P_Nh=\sum_{j=1}^N
c^{(N)}_j\ell_j\) are obtained from the finite Gram equations
\[
\left\langle\sum_{j=1}^N c^{(N)}_j\ell_j,\ell_i\right\rangle_{\Hl}
=\langle h,\ell_i\rangle_{\Hl}
=\langle h(\lambda_i),u_i\rangle_{E_*},
\qquad 1\leq i\leq N.
\]
If the Gram matrix is singular, its Moore--Penrose solution gives the
same projection.  Set
\begin{equation}
\label{eq:finite}
g_N(z)=\sum_{j=1}^N c^{(N)}_jK(z,(0,\lambda_j))u_j.
\end{equation}
The identity theorem and totality of the coefficient vectors imply that
the \(\ell_j\)'s have dense span in \(\Hl\).  Consequently,
\[
g_N=R^*P_Nh\longrightarrow R^*h=g
\quad\hbox{in }\Hk.
\]
Bounded point evaluations give pointwise convergence, and local boundedness
of \(K(z,z)\) gives uniform convergence on compact subsets of \(\D^2\).
Every datum on the right side of the finite Gram equations is a value of
the one-variable trace \(h\).

\section{The promised colligation-block formulas}

Let \(T=R^*:\Hl\to\Hk\) be the reconstruction operator in
Theorem~\ref{thm:reconstruct}.  In the notation of the source paper, for
\(f=(f_1,f_2)\) define the analytic one-variable functions
\[
q_f(\zeta)=
\begin{cases}
\dfrac{f_1(0,\zeta)-f_1(0)+f_2(0,\zeta)-f_2(0)}{\zeta},
&\zeta\ne0,\\[6pt]
\partial_2f_1(0)+\partial_2f_2(0),&\zeta=0,
\end{cases}
\]
and, for \(\eta\in E\),
\[
q_\eta(\zeta)=
\begin{cases}
\dfrac{\Phi(0,\zeta)-\Phi(0)}{\zeta}\eta,&\zeta\ne0,\\[6pt]
\partial_2\Phi(0)\eta,&\zeta=0.
\end{cases}
\]
The source paper proves that these traces belong to \(\Hl\).  Its
previously implicit unique functions are therefore
\begin{equation}
\label{eq:AB}
(Af)_2=Tq_f,\qquad (B\eta)_2=Tq_\eta.
\end{equation}
Their first components are then the source paper's structured Gleason
quotients
\[
(Af)_1(z)=
\frac{f_1(z)-f_1(0)+f_2(z)-f_2(0)-z_2(Tq_f)(z)}{z_1},
\]
\[
(B\eta)_1(z)=
\frac{(\Phi(z)-\Phi(0))\eta-z_2(Tq_\eta)(z)}{z_1},
\]
with the removable values at \(z_1=0\) filled in analytically.
Combining \eqref{eq:AB} with \eqref{eq:cross} or the Gram scheme
\eqref{eq:finite} gives the requested formulas using only vertical traces.

\section{Scope and novelty status}

This answers Question 5.6 for every coefficient Hilbert space covered by
the source paper.  It is not a new general RKHS theorem: the restriction
mechanism is classical.  The contribution of this note is the direct
identification of the paper's vertical uniqueness lemma with the injective
case of that theorem, plus the explicit cross-kernel and finite-Gram
reconstruction formulas.

A bounded search covered the run indexes; arXiv:1312.0177; the exact phrase
of Question 5.6; the terms \emph{canonical Agler}, \emph{vertical slice},
\emph{restriction}, and \emph{reconstruction}; and later arXiv papers on
Agler spaces and transfer-function realizations.  No later paper explicitly
answering Question 5.6 was found.  Because the decisive abstract theorem is
classical, the conservative provenance label is
\emph{literature-implied answer}, not a new full theorem.  The official AMS
metadata identifies Aronszajn's Section 5 as the restriction theorem.
An environment download limit prevented copying that supporting PDF into
this packet; the source PDF and its question screenshot are included.

\section{Review recommendation}

Verify the convention for operator-valued kernel inner products in
\eqref{eq:cross}, the source paper's vertical uniqueness argument immediately
before Question 5.6, and that the trace data \(q_f,q_\eta\) are the ones in
equations (5.8)--(5.9).  The reconstruction theorem itself is a two-line
adjoint argument once those points are fixed.

\begin{thebibliography}{9}
\bibitem{BickelKnese}
K. Bickel and G. Knese,
\emph{Canonical Agler Decompositions and Transfer Function Realizations},
Trans. Amer. Math. Soc. 368 (2016), 6293--6324;
arXiv:1312.0177.

\bibitem{Aronszajn}
N. Aronszajn,
\emph{Theory of reproducing kernels},
Trans. Amer. Math. Soc. 68 (1950), 337--404,
Section 5, doi:10.1090/S0002-9947-1950-0051437-7.
\end{thebibliography}

\end{document}
